% =============================================================
%  IEEE Conference Paper
%  Title : Continual Learning-Based User Embedding System
%          for Evolving Interest Representation
%  Style : IEEEtran (conference)
%  Notes : References forced to a fresh page in single-column.
%          Compile twice with pdflatex to resolve cross-refs.
% =============================================================

\documentclass[conference]{IEEEtran}

% ----- packages ----------------------------------------------
\usepackage[T1]{fontenc}
\usepackage[utf8]{inputenc}
\usepackage{amsmath}
\usepackage{amssymb}
\usepackage{graphicx}
\usepackage{booktabs}
\usepackage{array}
\usepackage{url}
\usepackage{cite}
\usepackage{algorithm}
\usepackage{algorithmic}
\usepackage{float}
\usepackage{microtype}
\usepackage[hidelinks]{hyperref}

% ----- notation shortcuts ------------------------------------
\newcommand{\eb}{\mathbf{e}}
\newcommand{\lam}{\lambda}
\newcommand{\alph}{\alpha}
\newcommand{\Fc}{\mathcal{F}}

% ----- minor spacing tweak -----------------------------------
\setlength{\parskip}{1.5pt}

% =============================================================
\begin{document}

% ----- title block -------------------------------------------
\title{%
  Continual Learning-Based User Embedding System\\[3pt]
  for Evolving Interest Representation%
}

\maketitle
\thispagestyle{empty}
\pagestyle{plain}

% =============================================================
%  ABSTRACT
% =============================================================
\begin{abstract}

Most recommender systems look at a user's profile as something
pretty static --- a vector that gets calculated once and only
updated when there is enough new data to bother retraining.
That is fine for users whose interests do not shift much, but
as soon as someone's tastes start changing, this approach
seriously falls apart. A learner who starts with Python basics,
moves into machine learning, and then dives deep into
transformer architectures is practically a different person
at each stage. If a system keeps referencing their early
interests months later, it is bound to recommend things they
no longer care about.

This paper digs into that problem with a four-layer pipeline.
The first layer builds a controlled dataset of 50 users
distributed across five technical domains, each following a
topic timeline that drifts gradually from beginner to advanced
material. The second layer converts those timelines into
384-dimensional semantic vectors using a pre-trained sentence
transformer --- no domain-specific fine-tuning required.
The third layer tests three different update strategies in
sequence --- a plain overwrite, Elastic Weight Consolidation
(EWC), and experience replay --- and measures how much of each
user's original representation survives after four update
steps. The fourth layer places the entire user population into
a shared two-dimensional space for trajectory analysis and
nearest-neighbour retrieval.

On the baseline evaluation, same-domain users showed an
average cosine similarity of $0.733$ compared to $0.221$
across domains --- a $3.32\times$ separation achieved with
no supervision whatsoever. Under plain overwriting, users
drifted a mean of $0.853$ from their initial state by the
final timestep, establishing a concrete forgetting baseline
that EWC and replay will be measured against in Phase~2.

\end{abstract}

\begin{IEEEkeywords}
continual learning,
user embeddings,
catastrophic forgetting,
representation learning,
sentence transformers,
elastic weight consolidation,
experience replay,
UMAP
\end{IEEEkeywords}

% =============================================================
%  I. INTRODUCTION
% =============================================================
\section{Introduction}

\IEEEPARstart{P}{icture} a student who starts the year
learning Python, drifts into machine learning, and by
year's end is knee-deep in transformer models. Their
interests are sharp and genuine at every point, but those
interests keep changing. If a system only remembers where
they began, it will keep pushing beginner material well
after they have moved on. And if it loses sight of their
starting point completely, it risks forgetting the key
foundations that still shape the way they think. Handling
this smoothly is not as easy as it sounds.

The standard engineering answer is periodic retraining:
gather the newest interaction data, rebuild the user model
on a schedule, and hope the update cadence stays fast
enough to remain relevant. That works well enough when
users are stable and compute is cheap. It breaks down in
two common situations. The first is when a user's interests
shift faster than the retraining cycle can keep up. The
second is when retraining on a long history ends up
dominated by recent behaviour, quietly overwriting
everything the user cared about a year ago.

Continual learning tackles the underlying problem head-on.
Rather than rebuilding from scratch, it asks how a learned
representation can be updated incrementally without
destroying what was already there. The core difficulty is
catastrophic forgetting~\cite{mccloskey1989} --- the
empirically well-established tendency of neural models to
overwrite earlier configurations when trained on new data,
even when that new data is closely related to the old.

Here, each user's embedding is treated as a small model in
its own right, and we study what happens as it is updated
sequentially across four timesteps. Three strategies are
evaluated: a naive overwrite serving as the worst-case
control, EWC~\cite{kirkpatrick2017} which uses per-dimension
importance weights to slow down changes to the most
significant parts of the old embedding, and experience
replay~\cite{rolnick2019} which keeps a short memory of
recent states and blends them into each update. Running all
three on the same synthetic dataset makes the
stability-plasticity trade-off visible in a way that is
concrete and directly measurable.

\smallskip
\noindent The specific contributions are:

\begin{itemize}
  \setlength\itemsep{2pt}

  \item A synthetic user data generator with configurable
        drift speed and cross-domain noise, producing
        reproducible topic timelines that mimic realistic
        learning progressions.

  \item A semantic embedding pipeline built on a frozen
        sentence transformer, including a snapshot mechanism
        that preserves each user's initial state as an
        immutable forgetting reference.

  \item Three continual learning update strategies
        implemented under a single consistent forgetting
        metric, enabling direct numerical comparison between
        them.

  \item A social embedding space layer supporting trajectory
        visualisation, pairwise similarity retrieval, and
        next-topic prediction by neighbour aggregation.

\end{itemize}

% =============================================================
%  II. RELATED WORK
% =============================================================
\section{Related Work}

\subsection{Catastrophic Forgetting and Continual Learning}

The problem of neural networks forgetting earlier tasks when
trained on new ones was formally documented by McCloskey and
Cohen~\cite{mccloskey1989}. For years it was treated more as
a quirky footnote than a major engineering concern. As people
started building larger sequential learning systems, the need
to actually fix it grew.

Kirkpatrick et al.~\cite{kirkpatrick2017} introduced EWC,
which uses the Fisher Information Matrix to identify parameters
that were most important for previous tasks and adds a
quadratic penalty to resist changing them. The intuition is
clean: not all parameters matter equally, so protect the ones
that do. Later work extended this to an online
variant~\cite{schwarz2018} that does not require storing
separate Fisher estimates per task, and to structured Kronecker
approximations~\cite{ritter2018} that handle correlated
parameters more carefully than the diagonal version does.

A different family of methods keeps a buffer of past samples
and replays them during future training. Lopez-Paz and
Ranzato~\cite{lopezpaz2017} framed this with explicit gradient
constraints to prevent interference. Rolnick et
al.~\cite{rolnick2019} showed that plain replay --- no fancy
constraints, just mix in old samples --- performs surprisingly
well across a range of benchmarks. Parisi et
al.~\cite{parisi2019} provide a thorough survey, splitting the
field into regularisation-based, architecture-based, and
replay-based approaches.

\subsection{Semantic User Modelling}

Collaborative filtering~\cite{koren2009} builds user
representations from interaction patterns. The resulting
vectors are good at predicting what a user will click on but
carry no semantic information about what the user actually
knows or cares about. Content-based methods fix this but
traditionally required hand-crafted feature pipelines.

The transformer revolution changed the picture considerably.
BERT~\cite{devlin2019} showed that dense semantic
representations could be learned from raw text without
task-specific labels. Reimers and
Gurevych~\cite{reimers2019} adapted that architecture into
Sentence-BERT, optimised specifically to produce sentence
embeddings useful for similarity tasks. Distilled descendants
like all-MiniLM-L6-v2~\cite{wang2020} compressed the model
enough to run comfortably on a CPU while preserving most of
the semantic quality --- which is what made them practical
for this work.

\subsection{Dynamic User Models}

Session-based recommendation~\cite{hidasi2016} models
short-window interaction sequences. Knowledge
tracing~\cite{piech2015} estimates which concepts a student
has mastered, but it requires explicit correctness signals
that are not available in general learning interest data.

The gap this paper addresses is longer-horizon semantic
evolution driven purely by topic keywords, without interaction
counts, session boundaries, or correctness labels. Applying
continual learning update algorithms directly to this setting
has not, to our knowledge, been studied in prior work.

% =============================================================
%  III. SYSTEM ARCHITECTURE
% =============================================================
\section{System Architecture}

The system is built as four independent layers connected in a
straight pipeline. Each layer takes one kind of input and
produces one kind of output, and none of them needs to know
anything about the internals of its neighbours. Swapping the
embedding model in Layer~2 does not touch a single line of the
continual learning code in Layer~3 --- the two layers
communicate only through a NumPy array.

\subsection{Layer 1 --- Synthetic Timeline Generation}

No publicly available longitudinal dataset records individual
learning interest trajectories at keyword resolution, so the
dataset was generated synthetically with controlled properties.

At the heart of the generator is a \texttt{TOPIC\_GRAPH}: a
dictionary mapping five learning personas --- machine learning
(\texttt{ml}), robotics, web development (\texttt{web}), data
science, and cybersecurity (\texttt{security}) --- to one or
more ordered topic progressions. Each progression is a list
of 8--10 technical keywords moving from foundational to
advanced. The \texttt{ml} persona has three separate paths:
one ending in large language model topics, one in MLOps, and
one in computer vision. Users within the same persona are
similar but not identical, because they can follow different
paths.

A timeline is built by sliding a window of width $w = 3$
along the selected progression at one step per timestep:

\begin{equation}
  \mathcal{T}(u,\,t)
  \;=\;
  P\bigl[\,t \;:\; t + w\,\bigr]
  \label{eq:window}
\end{equation}

The one-step slide creates a two-topic overlap between
consecutive timesteps, so a user who held topics
$\{A, B, C\}$ at step $t$ holds $\{B, C, D\}$ at step
$t{+}1$. That overlap is what keeps the drift gradual rather
than abrupt. An earlier version of the generator used a step
size of two, which caused the window to overshoot the end of
shorter progressions and wrap back, producing identical topic
sets at steps three and four. Reducing the step to one
eliminated the problem entirely.

Cross-domain noise is added with probability $p_n = 0.10$
per timestep: one beginner-level keyword from a different
persona is appended to the current step, modelling the
occasional off-topic curiosity that real learners show. The
50 users are assigned to personas in round-robin order (ten
per persona) and the whole generation process is seeded at
42 for reproducibility.

\subsection{Layer 2 --- Semantic Embedding}

Before encoding, the keywords in each timestep's topic list
are sorted alphabetically and joined into a
space-separated sentence. Sorting matters because transformer
models are sensitive to word order, and the order of topics
within a timestep carries no information about a user's
interests --- it is an artefact of how the generator stacked
them. Without sorting, the same three topics in two different
orders can produce vectors separated by a non-trivial cosine
distance, injecting noise into the forgetting measurements
downstream.

The sorted sentence is passed to
\texttt{all-MiniLM-L6-v2}, a six-layer transformer distilled
from a larger teacher model, producing a 384-dimensional
vector normalised to unit length. No part of the model is
fine-tuned; it functions purely as a frozen feature extractor.
Because all output vectors are unit vectors, cosine similarity
between any two of them is just their dot product --- no
division required.

All 250 embeddings ($50\text{ users} \times 5\text{
timesteps}$) are collected into a master array of shape
$(N, T, D) = (50, 5, 384)$ and saved as a NumPy binary file.
Before any update strategy runs, the $t{=}0$ embedding of
every user is saved separately as an immutable snapshot.
These snapshots are the forgetting reference throughout all
experiments and are never modified.

\subsection{Layer 3 --- Continual Learning Strategies}

\textbf{Forgetting metric.}\enspace
After every update step, the degree of forgetting for user
$u$ at timestep $t$ is:

\begin{equation}
  \Fc(u,\,t)
  \;=\;
  1 \;-\; \eb_{u,t}^{\top}\,\eb_{u,0}
  \label{eq:forgetting}
\end{equation}

where $\eb_{u,t}$ is the embedding after the update and
$\eb_{u,0}$ is the frozen initial snapshot. Both are unit
vectors, so the dot product is the cosine similarity. A score
of zero means nothing has changed; a score near one means the
embedding has moved about as far from its starting point as
it can. The mean across users,
$\overline{\Fc}(t) = N^{-1}\sum_u \Fc(u,t)$, gives one
point on the forgetting curve at each timestep.

\subsubsection{Baseline --- Sequential Overwrite}

The baseline replaces the current embedding with the new one,
unconditionally:

\begin{equation}
  \eb_{u,t} \;\leftarrow\; \eb_{\mathrm{new}}
  \label{eq:baseline}
\end{equation}

Every bit of historical information is discarded at every
step. This is the worst case for catastrophic forgetting and
serves as the control condition that EWC and replay are
measured against.

\subsubsection{Elastic Weight Consolidation}

EWC asks which of the 384 dimensions mattered most to the
previous embedding and applies resistance proportional to
that importance. The diagonal of the Fisher Information Matrix
is approximated by squaring the current embedding values ---
dimensions with large absolute values are assumed to carry
stronger signal and are therefore protected more aggressively:

\begin{equation}
  F_i
  \;=\;
  \bigl(\eb_{u,t-1,\,i}\bigr)^{2}
  \label{eq:fisher}
\end{equation}

The update for each dimension $i$ is then a weighted
interpolation between the old and new values:

\begin{equation}
  \eb_{u,t,\,i}
  \;=\;
  \frac{%
    \eb_{\mathrm{new},\,i}
    \;+\; \lam\,F_i\;\eb_{u,t-1,\,i}
  }{%
    1 \;+\; \lam\,F_i
  }
  \label{eq:ewc}
\end{equation}

The result is re-normalised to unit length. The
regularisation strength $\lam$ controls the trade-off: a
small value lets the embedding adapt quickly but forgets
more; a large value holds the old representation tightly but
adapts slowly. We test
$\lam \in \{0.1,\, 0.5,\, 1.0,\, 5.0\}$ and select the
value that minimises mean forgetting at the final step.

\subsubsection{Experience Replay}

The replay strategy sidesteps the importance-weight
calculation entirely. It keeps a first-in-first-out buffer
$\mathcal{B}_u$ of recent embeddings for each user, with
capacity $k$. Before every update, the current embedding is
pushed into the buffer (the oldest entry is evicted if the
buffer is full). The new embedding is then blended with the
buffer mean:

\begin{equation}
  \eb_{u,t}
  \;=\;
  (1 - \alph)\;\eb_{\mathrm{new}}
  \;+\;
  \frac{\alph}{|\mathcal{B}_u|}
  \sum_{\mathbf{b}\,\in\,\mathcal{B}_u} \mathbf{b}
  \label{eq:replay}
\end{equation}

followed by re-normalisation. The mixing ratio $\alph$
determines how strongly the memory influences each update.
Algorithm~\ref{alg:replay} gives the full procedure. We
evaluate $k \in \{1,\,2,\,3\}$ and
$\alph \in \{0.2,\,0.3,\,0.5\}$.

\begin{algorithm}[t]
  \caption{Experience Replay Embedding Update}
  \label{alg:replay}
  \begin{algorithmic}[1]
    \REQUIRE
      $\eb_{\mathrm{old}}$,\ $\eb_{\mathrm{new}}$,\
      buffer $\mathcal{B}$,\ capacity $k$,\ ratio $\alph$
    \STATE Push $\eb_{\mathrm{old}}$ onto $\mathcal{B}$
    \IF{$|\mathcal{B}| > k$}
      \STATE Evict oldest entry from $\mathcal{B}$
    \ENDIF
    \STATE
      $\mathbf{m}
      \,\leftarrow\,
      |\mathcal{B}|^{-1}
      \sum_{\mathbf{b}\,\in\,\mathcal{B}} \mathbf{b}$
    \STATE
      $\eb
      \,\leftarrow\,
      (1-\alph)\,\eb_{\mathrm{new}}
      + \alph\,\mathbf{m}$
    \RETURN $\eb\;/\;\|\eb\|_{2}$
  \end{algorithmic}
\end{algorithm}

\subsection{Layer 4 --- Social Embedding Space}

The final layer exposes three capabilities over the full user
population.

\textbf{Trajectory visualisation.}\enspace
All 250 embeddings are stacked into one matrix and reduced
jointly --- first by PCA to 50 dimensions, then by
UMAP~\cite{mcinnes2018} to two (\texttt{n\_neighbors}\,=\,8,
\texttt{min\_dist}\,=\,0.2, cosine metric). Projecting every
timestep together places them in a single shared coordinate
system, so an arrow drawn from user $u$'s $t{=}0$ position
to their $t{=}4$ position is a geometrically honest
representation of how far and in which direction they drifted.

\textbf{Pairwise similarity.}\enspace
The full $N \times N$ cosine similarity matrix is computed
as $\mathbf{S} = \mathbf{E}\mathbf{E}^{\top}$, where each
row of $\mathbf{E}$ is a unit-normalised user embedding.
Cluster quality is quantified with the silhouette
score~\cite{rousseeuw1987} under cosine distance.

\textbf{Next-topic prediction.}\enspace
For a query user, the $k$ nearest neighbours are retrieved
from $\mathbf{S}$. Topics appearing in those neighbours'
subsequent timesteps are frequency-counted, and the most
common topics the query user has not yet encountered are
returned as recommendations.

% =============================================================
%  IV. EXPERIMENTAL SETUP
% =============================================================
\section{Experimental Setup}

\subsection{Dataset}

The 50-user synthetic dataset contains ten users per persona,
each with five timesteps of three topics. Every experiment
uses \texttt{seed\,=\,42}; the generation code and the
resulting JSON file are included in the repository.

\subsection{Embedding Model}

All embeddings come from \texttt{all-MiniLM-L6-v2} via the
\texttt{sentence-transformers} library (version~4.1.0). The
model has six transformer layers and produces 384-dimensional
output. Inference ran entirely on CPU; encoding all 250
sentences takes roughly two seconds on a standard laptop.
No gradients were computed or applied to the model at any
point during the experiments.

\subsection{Evaluation}

The primary output is the mean forgetting curve
$\overline{\Fc}(t)$ for each strategy over
$t \in \{1,2,3,4\}$, anchored at
$\overline{\Fc}(0) = 0$ by definition. Secondary outputs
include per-persona cosine similarity statistics, the overall
separation ratio, and the silhouette score under cosine
distance. Table~\ref{tab:hyper} lists every hyperparameter
and the values tested.

\begin{table}[t]
  \caption{Hyperparameter Search Space}
  \label{tab:hyper}
  \centering
  \renewcommand{\arraystretch}{1.3}
  \begin{tabular}{lc}
    \toprule
    \textbf{Parameter}             & \textbf{Values tested}       \\
    \midrule
    EWC regularisation $\lam$      & $\{0.1,\ 0.5,\ 1.0,\ 5.0\}$ \\
    Replay capacity $k$            & $\{1,\ 2,\ 3\}$              \\
    Replay mixing ratio $\alph$    & $\{0.2,\ 0.3,\ 0.5\}$       \\
    UMAP \texttt{n\_neighbors}     & 8                             \\
    UMAP \texttt{min\_dist}        & 0.2                           \\
    Random seed (all experiments)  & 42                            \\
    \bottomrule
  \end{tabular}
\end{table}

% =============================================================
%  V. RESULTS
% =============================================================
\section{Results}

\subsection{Embedding Space Validation}

Before any continual learning experiment makes sense, the
embedding space needs to actually encode something meaningful.
If users from the same domain are no more similar to each
other than to users from completely different fields, the
whole exercise falls apart before it starts.

Table~\ref{tab:sim} shows that is not the case. The overall
separation ratio comes out at $3.32\times$ --- same-domain
users average $0.734$ similarity while cross-domain pairs
average $0.221$ --- and this was achieved without a single
labelled training example. The pre-trained model's existing
knowledge of technical English is doing all the work.
Keywords like \textit{ros}, \textit{kinematics}, and
\textit{path\_planning} already cluster in its representation
space because they appear together constantly in robotics
literature.

\begin{table}[t]
  \caption{Per-Persona Cosine Similarity at $t = 0$}
  \label{tab:sim}
  \centering
  \renewcommand{\arraystretch}{1.3}
  \begin{tabular}{lccc}
    \toprule
    \textbf{Persona}      & \textbf{Intra} & \textbf{Inter}
                          & \textbf{Separation} \\
    \midrule
    \texttt{ml}           & 0.782 & 0.203 & $+0.579$ \\
    \texttt{robotics}     & 0.765 & 0.199 & $+0.567$ \\
    \texttt{web}          & 0.891 & 0.034 & $+0.857$ \\
    \texttt{data science} & 0.744 & 0.288 & $+0.457$ \\
    \texttt{security}     & 0.711 & 0.200 & $+0.511$ \\
    \midrule
    \textbf{Overall}      & \textbf{0.734} & \textbf{0.221}
                          & $\mathbf{+0.513}$ \\
    \bottomrule
  \end{tabular}
\end{table}

The web persona stands furthest apart from everything else
(inter-similarity of just $0.034$). This makes sense: HTML,
CSS, React, and GraphQL share essentially no vocabulary with
the Python-heavy terminology of the other four domains. Data
science and machine learning, by contrast, share enough
common ground --- \textit{pandas}, \textit{sklearn},
\textit{statistics} --- that their inter-persona similarity
is the highest of any pair at around $0.50$.

Figure~\ref{fig:pca} puts this into visual form. The web
cluster sits cleanly away from the others; robotics and
machine learning share a region of space but remain
distinguishable. The pairwise similarity matrix
(Fig.~\ref{fig:matrix}) makes the block structure explicit:
bright diagonal squares confirm within-group cohesion, and
the notably dark row and column for web confirm just how
semantically isolated it is from the rest.

\begin{figure}[t]
  \centering
  \includegraphics[width=\columnwidth]{figures/fig1_pca_t0.png}
  \caption{%
    PCA projection of $t{=}0$ embeddings. Each point is one
    user; colour encodes persona. Web (green) sits clearly
    apart. The overlap between machine learning and data
    science reflects genuine shared vocabulary rather than
    a flaw in the embedding.%
  }
  \label{fig:pca}
\end{figure}

\begin{figure}[t]
  \centering
  \includegraphics[width=\columnwidth]{figures/fig3_similarity_matrix.png}
  \caption{%
    Pairwise cosine similarity matrix for all 50 users at
    $t{=}0$, sorted by persona. The five bright diagonal
    blocks confirm within-group cohesion. The dark row and
    column for web confirm its semantic isolation from every
    other domain.%
  }
  \label{fig:matrix}
\end{figure}

\subsection{Baseline Forgetting Measurements}

Table~\ref{tab:drift} shows how far each persona drifts from
its initial state by $t{=}4$ under plain overwriting. Web
drifts the most, reaching a mean of $0.958$. The reason is
the same one that made web so distinct in the similarity
analysis: the vocabulary gap between beginner and advanced
topics is steep. A user whose initial topics were
\textit{html} and \textit{css} ends up at \textit{graphql}
and \textit{prisma} --- words that barely overlap in the
model's space even though the user has simply gone deeper
into web development rather than switching fields. Robotics
drifts the least ($0.777$), likely because the domain's
vocabulary --- sensors, control, kinematics, path planning
--- stays technically consistent from entry level to expert
level.

\begin{table}[t]
  \caption{Baseline Forgetting at $t = 4$ by Persona}
  \label{tab:drift}
  \centering
  \renewcommand{\arraystretch}{1.3}
  \begin{tabular}{lcc}
    \toprule
    \textbf{Persona}      & \textbf{Mean $\overline{\Fc}$}
                          & \textbf{Range}      \\
    \midrule
    \texttt{ml}           & 0.850 & $[0.53,\;1.06]$ \\
    \texttt{robotics}     & 0.777 & $[0.67,\;0.82]$ \\
    \texttt{web}          & 0.958 & $[0.93,\;1.00]$ \\
    \texttt{data science} & 0.798 & $[0.57,\;0.91]$ \\
    \texttt{security}     & 0.881 & $[0.84,\;0.95]$ \\
    \midrule
    \textbf{All users}    & \textbf{0.853}
                          & $[0.53,\;1.06]$     \\
    \bottomrule
  \end{tabular}
\end{table}

Figure~\ref{fig:traj} shows the same pattern spatially.
Web users converge toward a shared endpoint at the bottom
of the chart; machine learning arrows are shorter and point
in different directions depending on which progression path
each user followed. Figure~\ref{fig:drift} shows the full
forgetting curve over time. Every persona follows the same
rough shape: a steep initial jump at $t{=}1$, then a gradual
flattening as the embedding settles into its new region of
the space. The steepest early rise belongs to web; the
shallowest overall trajectory belongs to robotics.

\begin{figure}[t]
  \centering
  \includegraphics[width=\columnwidth]{figures/fig2_trajectories.png}
  \caption{%
    User trajectories from $t{=}0$ (faded dot) to $t{=}4$
    (solid dot) in PCA space. Arrow length corresponds to
    total semantic drift. Web users (green) converge on a
    shared advanced-topic cluster. Machine learning arrows
    are shorter and more dispersed, reflecting multiple
    available progression paths.%
  }
  \label{fig:traj}
\end{figure}

\begin{figure}[t]
  \centering
  \includegraphics[width=\columnwidth]{figures/fig4_drift_over_time.png}
  \caption{%
    Baseline semantic drift $\overline{\Fc}(t)$ over time.
    Thin lines are individual users; thick lines are
    persona-level averages. This curve is the control
    condition for Phase~2: EWC and replay buffer results
    will be overlaid on this same chart.%
  }
  \label{fig:drift}
\end{figure}

\subsection{Continual Learning Evaluation (Phase~2)}

The EWC and replay buffer experiments are currently running.
Given the size of the baseline forgetting signal ($0.853$)
and the structured nature of the synthetic progressions,
two success thresholds are defined: a 25\% reduction for EWC
($\overline{\Fc}_{\mathrm{EWC}} < 0.640$) and 20\% for
replay ($\overline{\Fc}_{\mathrm{replay}} < 0.682$). These
are intentionally conservative relative to what EWC achieves
on structured sequential tasks in the
literature~\cite{kirkpatrick2017}.

The final comparison figure will overlay all three forgetting
curves on a single chart. A consistent gap between the
baseline line and both CL lines across $t \in \{2,3,4\}$,
supported by paired $t$-tests, will confirm the central claim
of the paper.

% =============================================================
%  VI. DISCUSSION
% =============================================================
\section{Discussion}

\subsection{What Phase~1 Tells Us}

The $3.32\times$ separation ratio is the most practically
significant result from Phase~1, and it comes with an
important qualifier: none of it required labelled data. The
model was not told that robotics users should cluster
together. It already knew that \textit{kinematics} and
\textit{slam} are related concepts, and that \textit{html}
is not, because it saw those relationships played out across
a vast amount of technical text during pre-training.

This has a useful implication for deployment. In many real
systems, user interaction data is sparse or unavailable
early on --- the cold-start problem. A topic-keyword approach
built on a general-purpose sentence transformer could provide
a workable first approximation of user interests without
requiring any interaction history at all. Whether that
approximation is accurate enough to be useful in practice is
an empirical question, but the Phase~1 results suggest it is
at least worth investigating seriously.

The high baseline drift values (mean $0.853$) also deserve
a closer look than they might initially receive. Part of that
number is genuine interest evolution --- a user moving from
Python basics to large language models really has changed
significantly. But another part is a measurement artefact.
When \textit{python} and \textit{numpy} give way to
\textit{transformers} and \textit{llms}, the cosine distance
between those keyword sets is large in the model's space even
though the user has simply gone deeper into machine learning
rather than leaving it. A model fine-tuned on domain-specific
text would likely compress that distance, producing lower
apparent drift for the same actual learning progression.
That distinction matters when interpreting the Phase~2
forgetting reductions.

\subsection{Limitations}

The dataset is synthetic, and that is the most significant
caveat. The sliding-window mechanism produces smooth,
predictable drift that real learner behaviour almost
certainly does not exhibit. People revisit old topics,
abandon subjects for months, and pick up parallel interests
in unrelated areas. A system built and validated only on
smooth progressions might handle abrupt jumps poorly. The
most important next step is testing against a real learning
log dataset --- even a small one would be informative.

The second limitation is that all keywords are weighted
equally. A user who has spent three months working through
every tutorial on a topic and someone who glanced at it once
produce the same embedding update if their keyword lists look
the same. Incorporating an engagement signal --- even
something as simple as dwell time on a page --- would make
the updates meaningfully more informative.

The third is the diagonal Fisher approximation in EWC. If the
information encoding a user's core interest is distributed
across a cluster of correlated dimensions rather than
concentrated in a few individually large ones, the diagonal
approximation will underprotect it. Structured
approximations~\cite{ritter2018} handle this case better,
though at a higher computational cost.

\subsection{Where This Goes Next}

Once Phase~2 results are in, the natural next question is
cluster quality: does EWC preserve the within-persona
grouping better than plain overwriting? If the silhouette
score under EWC-updated embeddings stays closer to its
$t{=}0$ value than the baseline silhouette does, that is a
second piece of evidence that EWC is actually protecting
meaningful structure rather than just slowing down change.
The trajectory visualisation will give an intuitive view of
the same thing: EWC arrows should be shorter on average than
baseline arrows, and they should point in directions that
stay within the user's original cluster rather than crossing
into neighbouring ones.

Beyond Phase~2, the modular design means replacing the
synthetic dataset with a real one requires changes only in
Layer~1. Every other layer operates on the same $(N, T, D)$
array regardless of how the topic lists were obtained. That
separation was a deliberate choice during implementation, and
it should make the transition to real data considerably less
painful.

% =============================================================
%  VII. CONCLUSION
% =============================================================
\section{Conclusion}

This paper set out to answer a practical question: can
simple, computationally cheap continual learning strategies
meaningfully reduce catastrophic forgetting in a sequential
user embedding pipeline? Building the infrastructure to
answer that question rigorously turned out to be most of the
work, and that infrastructure is now in place.

Phase~1 established two things. The embedding space built
from topic keywords and a frozen sentence transformer is
genuinely well-organised by domain, achieving a $3.32\times$
same-domain versus cross-domain similarity ratio without any
supervision. And naive sequential overwriting erases a
substantial fraction of each user's initial representation
--- mean drift of $0.853$ after four steps --- providing a
strong baseline signal for Phase~2 to improve upon.

If EWC and replay buffer each reduce that forgetting by
20--25\% in the Phase~2 evaluation, the result carries a
straightforward practical message: you do not need an
elaborate lifelong learning system to do meaningfully better
than naive overwriting in this setting. Simple strategies,
applied consistently and measured carefully, are enough to
preserve a noticeable portion of each user's identity across
sequential updates. That is a useful thing to know for anyone
building systems that need to track how people's interests
evolve over time.

% =============================================================
%  REFERENCES  ---  new page, single column
% =============================================================
\newpage
\onecolumn

\section*{References}

\vspace{4pt}

\begin{thebibliography}{15}

\bibitem{kirkpatrick2017}
J.~Kirkpatrick, R.~Pascanu, N.~Rabinowitz, J.~Veness,
G.~Desjardins, A.~A.~Rusu, K.~Milan, J.~Quan, T.~Ramalho,
A.~Grabska-Barwinska, D.~Hassabis, C.~Clopath, D.~Kumaran,
and R.~Hadsell,
``Overcoming catastrophic forgetting in neural networks,''
\textit{Proceedings of the National Academy of Sciences},
vol.~114, no.~13, pp.~3521--3526, 2017.

\bibitem{rolnick2019}
D.~Rolnick, A.~Ahuja, J.~Schwarz, T.~P.~Lillicrap,
and G.~Wayne,
``Experience replay for continual learning,''
in \textit{Advances in Neural Information Processing Systems
(NeurIPS)}, 2019.

\bibitem{mccloskey1989}
M.~McCloskey and N.~J.~Cohen,
``Catastrophic interference in connectionist networks:
The sequential learning problem,''
\textit{Psychology of Learning and Motivation},
vol.~24, pp.~109--165, 1989.

\bibitem{lopezpaz2017}
D.~Lopez-Paz and M.~Ranzato,
``Gradient episodic memory for continual learning,''
in \textit{Advances in Neural Information Processing Systems
(NeurIPS)}, 2017.

\bibitem{parisi2019}
G.~I.~Parisi, R.~Kemker, J.~L.~Part, C.~Kanan,
and S.~Wermter,
``Continual lifelong learning with neural networks:
A review,''
\textit{Neural Networks}, vol.~113, pp.~54--71, 2019.

\bibitem{schwarz2018}
J.~Schwarz, J.~Luketina, W.~M.~Czarnecki, A.~Grabska-Barwinska,
Y.~W.~Teh, R.~Pascanu, and R.~Hadsell,
``Progress \& compress: A scalable framework for continual
learning,''
in \textit{Proceedings of the International Conference on
Machine Learning (ICML)}, 2018.

\bibitem{ritter2018}
H.~Ritter, A.~Botev, and D.~Barber,
``Online structured Laplace approximations for overcoming
catastrophic forgetting,''
in \textit{Advances in Neural Information Processing Systems
(NeurIPS)}, 2018.

\bibitem{koren2009}
Y.~Koren, R.~Bell, and C.~Volinsky,
``Matrix factorization techniques for recommender systems,''
\textit{Computer}, vol.~42, no.~8, pp.~30--37, 2009.

\bibitem{devlin2019}
J.~Devlin, M.-W.~Chang, K.~Lee, and K.~Toutanova,
``BERT: Pre-training of deep bidirectional transformers for
language understanding,''
in \textit{Proceedings of NAACL-HLT}, 2019.

\bibitem{reimers2019}
N.~Reimers and I.~Gurevych,
``Sentence-BERT: Sentence embeddings using Siamese
BERT-networks,''
in \textit{Proceedings of EMNLP-IJCNLP}, 2019.

\bibitem{wang2020}
W.~Wang, F.~Wei, L.~Dong, H.~Bao, N.~Yang,
and M.~Zhou,
``MiniLM: Deep self-attention distillation for task-agnostic
compression of pre-trained transformers,''
in \textit{Advances in Neural Information Processing Systems
(NeurIPS)}, 2020.

\bibitem{hidasi2016}
B.~Hidasi, A.~Karatzoglou, L.~Baltrunas, and D.~Tikk,
``Session-based recommendations with recurrent neural
networks,''
in \textit{Proceedings of the International Conference on
Learning Representations (ICLR)}, 2016.

\bibitem{piech2015}
C.~Piech, J.~Bassen, J.~Huang, S.~Ganguli, M.~Sahami,
L.~J.~Guibas, and J.~Sohl-Dickstein,
``Deep knowledge tracing,''
in \textit{Advances in Neural Information Processing Systems
(NeurIPS)}, 2015.

\bibitem{mcinnes2018}
L.~McInnes, J.~Healy, and J.~Melville,
``UMAP: Uniform manifold approximation and projection for
dimension reduction,''
\textit{arXiv preprint arXiv:1802.03426}, 2018.

\bibitem{rousseeuw1987}
P.~J.~Rousseeuw,
``Silhouettes: A graphical aid to the interpretation and
validation of cluster analysis,''
\textit{Journal of Computational and Applied Mathematics},
vol.~20, pp.~53--65, 1987.

\end{thebibliography}

\end{document}
